\section*{Introduction}

Tout en se montrant critique de l'évolution des formes du pouvoir terrestre, Ambroise de Milan se résout à accepter la monarchie, et même à composer avec au travers de sa réflexion politique. Cette acceptation est principalement liée à sa volonté de s'impliquer dans la vie politique contemporaine. L'objectif de l'évêque, au-delà de son rôle auprès de son clergé, est de se maintenir au plus près du pouvoir pour faciliter son influence. Avant un troisième chapitre qui se concentrera sur les actions concrètes de l'évêque face aux empereurs, ce deuxième chapitre vise à rentrer plus en profondeur dans sa théorie politique au sujet de la relation entre foi chrétienne et exercice du pouvoir. Il n'est plus question d'observer la théorie politique abstraite d'Ambroise, mais d'analyser la vision ambrosienne du comportement et des actions attendus des empereurs. L'enjeu est de comprendre comment l'évêque cherche à modeler les souverains romains par l'exigence morale et la nécessité d'incarner et de défendre les vertus chrétiennes.

\bigskip

Ambroise ne pense le monde qui l'entoure, et donc l'exercice du pouvoir puisqu'il y est constamment confronté, qu'à travers le prisme du christianisme. C'est dans ce cadre que s'inscrit sa vision du pouvoir : puisque le pouvoir vient de Dieu, il est naturel qu'il finisse dans les mains d'un chrétien, et si ce pouvoir est chrétien, alors il doit être tout aussi naturel qu'il ne s'exerce que dans l'objectif de protéger la foi chrétienne. Le bon souverain imaginé par Ambroise est directement lié à la lutte contre les autres courants religieux. C'est le cœur de ce chapitre : la foi de l'empereur impacte la société, le règne de l'empereur mais aussi la réussite de sa descendance. Cette logique explique notamment la part importante qu'occupe la valorisation de la lutte contre le polythéisme dans son œuvre. C'est ce travail qui amène certains historiens à montrer l'influence d'Ambroise dans les décisions impériales de Gratien ou Théodose entre 380 et 392 contre les cultes païens.

\bigskip

Sa réflexion pour modeler le monarque idéal ne part évidemment pas de rien. Il s'appuie sur des sources d'autorité reconnues par son auditoire, et donc par l'aristocratie romaine chrétienne. C'est ainsi que, dans sa construction de l'empereur romain, apparaît forcément l'influence biblique. Les figures royales issues des livres de Samuel ou des Rois sont régulièrement mobilisées pour servir de modèles de conduite face à la morale et au péché, ou alors de contre-exemples de souverains ayant échoué à maintenir la concorde de la société, souvent à cause d'un manque de soutien divin, et donc de foi. Une grande partie de sa réflexion se base sur les comparaisons entre les rois de l'Ancien Testament et les empereurs qu'il côtoie. Sa forte utilisation de la Bible doit également être reliée à sa connaissance de la philosophie classique, principalement cicéronienne. Il cherche à s'en détacher, mais reconnaît une proximité intellectuelle et l'utilise notamment pour forger un pouvoir idéal, encadré par la justice.

\bigskip

Et plus important encore pour ce chapitre, c'est vraiment la relation de l'empereur à la foi qui constitue le cœur de la réflexion politique d'Ambroise. Notre analyse s'inscrira dans cette perspective : le respect des volontés divines dans les actions terrestres, la valorisation de la piété pour la sauvegarde du règne et de la dynastie, et l'obéissance aux exigences chrétiennes mises en avant par Ambroise.

\bigskip

L'étude de la théorie du modelage des empereurs se fera de nouveau grandement à partir du \latin{De officiis}. Le traité offre une grande analyse du lien qu’a l’évêque avec la philosophie cicéronienne, ce qui permet déjà de mieux situer intellectuellement Ambroise. C'est également un ouvrage qui amène une meilleure compréhension des vertus qu’il veut valoriser et transposer aux souverains. Pour rentrer un peu plus dans la question de la pensée politique, le \latin{De obitu Theodosii} se révèle fondamental. L'oraison funèbre lui sert d'outil politique pour à la fois valoriser Théodose en dressant un bilan du monarque idéal, stabiliser la suite de l'Empire en assurant la légitimité du jeune Honorius et la confiance de l'armée à son propos, et surtout ancrer ses propos derrière les figures bibliques qui sont longuement invoquées. Enfin, ce chapitre utilisera certaines œuvres directement liées à la Bible comme son Apologie de David pour évoquer le lien entre le roi d'Israël et Théodose, ou son commentaire de l'Évangile selon Luc pour questionner l'origine du pouvoir.

\bigskip

Afin de comprendre comment l'évêque de Milan souhaite forger l'empereur chrétien, nous analyserons dans un premier temps l'ensemble des fondements intellectuels qui influencent sa pensée, en étudiant la manière dont il s'approprie les figures scripturaires et transforme les réflexions philosophiques en vertus chrétiennes. Puis nous démontrerons l'importance de la chrétienté dans l'exercice concret du pouvoir, et la façon dont Ambroise fait de cette exigence un outil d'influence sur les monarques.
